\scanpage{16}
Passant à la limite projective, on trouve
\[
\unsure{\mathcal G|\mathscr G}_\infty=\mathrm{Aut}_{\unsure{K|k}}(\underset{\substack{\vert\\\inserted{\substack{\text{rev. abélien}\\\text{étale maximal}\\\text{de }\mathscr X}}}}{\unsure{\mathscr X|\mathcal X|\mathfrak X}_\infty}\mathbin{-}\mathscr X)\simeq\Gamma\mathbin{\cdot}T_\infty(\overline{\unsure{\mathbb Q|Q}}^{\unsure{*|\times}})
\]

-- -- -

C) $\mathrm{Card}\,S=3\qquad\unsure{\mathfrak G|G|\mathscr G}\xrightarrow{\sim}\mathfrak G_S\quad(\simeq\mathfrak G_3)$

il y a une façon canonique de réduire le corps de base à $k$, de façon que $S$ devienne \unsure{rat.|ratl} $/k$. On est conduit ici à la tour de Fermat :
\[
\mathcal G_n=\bigl(\Gamma\times\deleted{\unsure{\pi}}\mathfrak G_S\bigr)\mathbin{\cdot}\underbrace{\pi_n}_{\mathfrak G_{(n)}}\qquad \pi_n=\mu_n^{S}/\mu_n
\]
\[
\deleted{\underbrace{\unsure{\Gamma\times\mathfrak G_S}}_{\unsure{\mathfrak G_{(n)}}}}
\]
\unsure{..---|...}

D) $\mathrm{Card}\,S\geq4$. Soit \deleted{\unsure{Card}} $s=\unsure{\{|(}s_0,s_1,s_\infty\unsure{\}|)}$ \inserted{un triplet de} trois points \unsure{distinct|distincts} de $S$, on les utilise pour définir un iso
\[
\widehat{\mathscr X}\underset{\varphi_s}{\simeq}\mathbb P^1_{\overline{\mathbb Q}}
\]
transformant $s_0$ en $0$, $s_1$ en $1$, $s_\infty$ en $\infty$. Les autres points \deleted{sont} \inserted{de $S$ s’identifient} donc à des \unsure{élé\textsuperscript{ts}|éléments} de $\overline{\mathbb Q}\setminus\{0,1\}$, soit $K_{\unsure{s|\xi}}\subset\overline{\mathbb Q}$
% archive sheet number: 15
